% Abstract approximately around 200 words

\begin{abstract}
Large language models can produce convincing explanations
without necessarily acknowledging factors associated with
their answers. Whether this behaviour persists in current
reasoning models remains an important question for
explainable AI. We investigate misleading-cue susceptibility
and visible explanation acknowledgement through a controlled
adaptation of an earlier Suggested Answer experiment.
GPT-5.6 Sol and Claude Opus 5.5 were evaluated on 270
BIG-Bench Hard questions spanning 13 tasks, under original,
misleading-cue, and uncued repeat conditions. Across 540
cued responses, GPT followed the suggested incorrect
answer in nine cases (3.3\%), compared with one case
(0.4\%) for Claude. The repeat condition revealed
additional answer variability without misleading cues.
Two independent human raters agreed that nine of the ten
cue-following explanations did not acknowledge the
misleading suggestion (Cohen's $\kappa = 1.00$).
Although the small number of cases and ordinary response
variability limit causal conclusions, the findings
highlight an important distinction: infrequent
cue-following does not guarantee that visible
explanations acknowledge potentially influential
suggestions. Evaluating both behavioural susceptibility
and explanation acknowledgement remains essential
when assessing explanation faithfulness.
\end{abstract}